\documentclass[11pt,letterpaper]{article}
\usepackage[technical-reference]{cornell-notes}
\usepackage{needspace}

\title{Clause 26: Ranges Library - Cornell Notes}
\author{}
\date{}
\setDocTitle{Clause 26: Ranges Library}
\setDocSubtitle{C++ 2024 Cornell Notes}
\setDocOwner{C++ 2024 Cornell Notes}
\setDocAuthor{}
\setDocDate{}
\setCornellCollection{C++ 2024 Cornell Notes}
\setCornellUnitType{Clause}
\setCornellUnitNumber{26}
\setCornellUnitTitle{Ranges Library}
\setCornellCompanionLabel{C++ Technical Reference}
\hypersetup{pdftitle={Clause 26: Ranges Library - Cornell Notes},pdfsubject={C++ technical study notes},pdfauthor={}}

\begin{document}
\maketitle
\begin{CornellReferenceOrientationBox}
\textbf{Purpose.} Defines range concepts, views, adaptors, factories, customization points, and machinery for composing lazy sequence transformations while tracking borrowing, commonality, sizing, and iterator validity.
\par\textbf{Role.} The composable sequence-view layer used by range-aware algorithms.
\par\textbf{Principal dependencies.} Uses the library-wide rules of Clause 16 together with the applicable core-language concepts and supporting library clauses.
\par\textbf{Primary audience.} generic programmers and view authors.
\par\textbf{Major subject groups.} General; Header <ranges> synopsis; Range access; Range requirements; Range utilities; Range factories; Range adaptors; Range generators.
\end{CornellReferenceOrientationBox}
\section{Learning Objectives}
\begin{itemize}[label=\textcolor{CornellReferenceTeal}{$\blacktriangleright$}]
\item Define the governing ideas of General and apply its stated contract at a call site.
\item Identify the governing ideas of Header <ranges> synopsis and apply its stated contract at a call site.
\item Distinguish the governing ideas of Range access and apply its stated contract at a call site.
\item Explain the governing ideas of Range requirements and apply its stated contract at a call site.
\item Trace the governing ideas of Range utilities and apply its stated contract at a call site.
\item Apply the governing ideas of Range factories and apply its stated contract at a call site.
\item Analyze the governing ideas of Range adaptors and apply its stated contract at a call site.
\item Evaluate the governing ideas of Range generators and apply its stated contract at a call site.
\item Diagnose the governing ideas of General and apply its stated contract at a call site.
\item Compare the governing ideas of ranges::begin and apply its stated contract at a call site.
\item Verify the governing ideas of ranges::end and apply its stated contract at a call site.
\item Relate the governing ideas of ranges::cbegin and apply its stated contract at a call site.
\item Classify the governing ideas of ranges::cend and apply its stated contract at a call site.
\item Justify the governing ideas of ranges::rbegin and apply its stated contract at a call site.
\item Audit the governing ideas of ranges::rend and apply its stated contract at a call site.
\item Summarize the governing ideas of ranges::crbegin and apply its stated contract at a call site.
\end{itemize}
\section{Normative Structure Map}
\small\begin{longtable}{@{}>{\RaggedRight\arraybackslash}p{0.13\textwidth}>{\RaggedRight\arraybackslash}p{0.27\textwidth}>{\RaggedRight\arraybackslash}p{0.19\textwidth}>{\RaggedRight\arraybackslash}p{0.31\textwidth}@{}}
\toprule\textbf{Subclause} & \textbf{Subject} & \textbf{Rule category} & \textbf{Key dependencies / entities}\\\midrule\endfirsthead
\toprule\textbf{Subclause} & \textbf{Subject} & \textbf{Rule category} & \textbf{Key dependencies / entities}\\\midrule\endhead
\bottomrule\endfoot
26.1 & General & library semantics & Uses the library-wide rules of Clause 16 together with the applicable core-language concepts and supporting library clauses.\\
26.2 & Header <ranges> synopsis & interface / declarations & Uses the library-wide rules of Clause 16 together with the applicable core-language concepts and supporting library clauses.\\
26.3 & Range access & library semantics & General, ranges::begin, ranges::end, ranges::cbegin, and related subclauses\\
26.4 & Range requirements & requirements & General, Ranges, Sized ranges, Views, and related subclauses\\
26.5 & Range utilities & library semantics & General, Helper concepts, View interface, Sub-ranges, and related subclauses\\
26.6 & Range factories & library semantics & General, Empty view, Single view, Iota view, and related subclauses\\
26.7 & Range adaptors & library semantics & General, Range adaptor objects, Movable wrapper, Non-propagating cache, and related subclauses\\
26.8 & Range generators & library semantics & Overview, Header <generator> synopsis, Class template generator, Members, and related subclauses\\
\end{longtable}\normalsize
\begin{CornellReferenceNormativeBox}A heading maps the clause; it does not replace the detailed conditions. For each operation, read requirements, effects, result, exceptions, complexity, remarks, and notes with their stated normative force.\end{CornellReferenceNormativeBox}
\subsection*{Rule Classification Guide}
\small\begin{longtable}{@{}>{\RaggedRight\arraybackslash}p{0.25\textwidth}>{\RaggedRight\arraybackslash}p{0.67\textwidth}@{}}\toprule\textbf{Label or category} & \textbf{How to read it}\\\midrule\endfirsthead\toprule\textbf{Label or category} & \textbf{How to read it (continued)}\\\midrule\endhead\bottomrule\endfoot
Constraints & Conditions controlling whether a declaration, overload, or specialization participates in the relevant context.\\
Mandates & Requirements checked for the described declaration or specialization; failure has the stated well-formedness consequence.\\
Preconditions & Obligations the caller establishes before evaluation; violating one forfeits the operation's contract.\\
Effects / Returns / Postconditions & Separate the state change, produced result, and state guaranteed after successful completion.\\
Throws / Error conditions & State how failure is reported and which exceptional paths are permitted or required.\\
Complexity & Bounds or exact counts for the operations named by the specification.\\
Remarks / Recommended practice & Remarks refine applicability or participation; recommended practice guides implementations without silently becoming a program requirement.\\
Exposition-only & A device for describing semantics, not a public name on which a program can depend.\\
\end{longtable}\normalsize
\section{Cornell Cue-and-Notes}
\Needspace{8\baselineskip}
\subsection*{26.1 General}
\begin{longtable}{@{}>{\RaggedRight\arraybackslash\columncolor{CornellReferenceCueGray}}p{0.28\textwidth}>{\RaggedRight\arraybackslash}p{0.64\textwidth}@{}}
\rowcolor{CornellReferenceNavy}\color{white}\textbf{Cue / recall prompt} & \color{white}\textbf{Notes}\\\endfirsthead
\rowcolor{CornellReferenceNavy}\color{white}\textbf{Cue / recall prompt} & \color{white}\textbf{Notes (continued)}\\\endhead
\arrayrulecolor{CornellReferenceLineGray}
\textbf{What contract governs General?}\\[-1mm]{\scriptsize\CornellClauseRef{26.1}} & Frames the terminology, applicability, and relationships needed to interpret General; detailed rules remain in the following subclauses.\\\hline
\end{longtable}
\begin{CornellReferenceSummaryBox}General supplies the library semantics portion of this clause. The key review move is to connect its local wording to the clause-wide model, then classify each condition by normative force before relying on it.\end{CornellReferenceSummaryBox}
\Needspace{8\baselineskip}
\subsection*{26.2 Header <ranges> synopsis}
\begin{longtable}{@{}>{\RaggedRight\arraybackslash\columncolor{CornellReferenceCueGray}}p{0.28\textwidth}>{\RaggedRight\arraybackslash}p{0.64\textwidth}@{}}
\rowcolor{CornellReferenceNavy}\color{white}\textbf{Cue / recall prompt} & \color{white}\textbf{Notes}\\\endfirsthead
\rowcolor{CornellReferenceNavy}\color{white}\textbf{Cue / recall prompt} & \color{white}\textbf{Notes (continued)}\\\endhead
\arrayrulecolor{CornellReferenceLineGray}
\textbf{What contract governs Header <ranges> synopsis?}\\[-1mm]{\scriptsize\CornellClauseRef{26.2}} & Catalogs the declarations made available by Header <ranges> synopsis. The synopsis is an interface map; the detailed specifications supply constraints and behavior.\\\hline
\end{longtable}
\begin{CornellReferenceSummaryBox}Header <ranges> synopsis supplies the interface / declarations portion of this clause. The key review move is to connect its local wording to the clause-wide model, then classify each condition by normative force before relying on it.\end{CornellReferenceSummaryBox}
\Needspace{8\baselineskip}
\subsection*{26.3 Range access}
\begin{longtable}{@{}>{\RaggedRight\arraybackslash\columncolor{CornellReferenceCueGray}}p{0.28\textwidth}>{\RaggedRight\arraybackslash}p{0.64\textwidth}@{}}
\rowcolor{CornellReferenceNavy}\color{white}\textbf{Cue / recall prompt} & \color{white}\textbf{Notes}\\\endfirsthead
\rowcolor{CornellReferenceNavy}\color{white}\textbf{Cue / recall prompt} & \color{white}\textbf{Notes (continued)}\\\endhead
\arrayrulecolor{CornellReferenceLineGray}
\textbf{What contract governs General?}\\[-1mm]{\scriptsize\CornellClauseRef{26.3.1}} & Frames the terminology, applicability, and relationships needed to interpret General; detailed rules remain in the following subclauses.\\\hline
\textbf{What contract governs ranges::begin?}\\[-1mm]{\scriptsize\CornellClauseRef{26.3.2}} & Explains the traversal or view contract for ranges::begin; prove domain validity, lifetime, sentinel compatibility, and any borrowing or invalidation property.\\\hline
\textbf{What contract governs ranges::end?}\\[-1mm]{\scriptsize\CornellClauseRef{26.3.3}} & Explains the traversal or view contract for ranges::end; prove domain validity, lifetime, sentinel compatibility, and any borrowing or invalidation property.\\\hline
\textbf{What contract governs ranges::cbegin?}\\[-1mm]{\scriptsize\CornellClauseRef{26.3.4}} & Explains the traversal or view contract for ranges::cbegin; prove domain validity, lifetime, sentinel compatibility, and any borrowing or invalidation property.\\\hline
\textbf{What contract governs ranges::cend?}\\[-1mm]{\scriptsize\CornellClauseRef{26.3.5}} & Explains the traversal or view contract for ranges::cend; prove domain validity, lifetime, sentinel compatibility, and any borrowing or invalidation property.\\\hline
\textbf{What contract governs ranges::rbegin?}\\[-1mm]{\scriptsize\CornellClauseRef{26.3.6}} & Explains the traversal or view contract for ranges::rbegin; prove domain validity, lifetime, sentinel compatibility, and any borrowing or invalidation property.\\\hline
\textbf{What contract governs ranges::rend?}\\[-1mm]{\scriptsize\CornellClauseRef{26.3.7}} & Explains the traversal or view contract for ranges::rend; prove domain validity, lifetime, sentinel compatibility, and any borrowing or invalidation property.\\\hline
\textbf{What contract governs ranges::crbegin?}\\[-1mm]{\scriptsize\CornellClauseRef{26.3.8}} & Explains the traversal or view contract for ranges::crbegin; prove domain validity, lifetime, sentinel compatibility, and any borrowing or invalidation property.\\\hline
\textbf{What contract governs ranges::crend?}\\[-1mm]{\scriptsize\CornellClauseRef{26.3.9}} & Explains the traversal or view contract for ranges::crend; prove domain validity, lifetime, sentinel compatibility, and any borrowing or invalidation property.\\\hline
\textbf{What contract governs ranges::size?}\\[-1mm]{\scriptsize\CornellClauseRef{26.3.10}} & Explains the traversal or view contract for ranges::size; prove domain validity, lifetime, sentinel compatibility, and any borrowing or invalidation property.\\\hline
\textbf{What contract governs ranges::ssize?}\\[-1mm]{\scriptsize\CornellClauseRef{26.3.11}} & Explains the traversal or view contract for ranges::ssize; prove domain validity, lifetime, sentinel compatibility, and any borrowing or invalidation property.\\\hline
\textbf{What contract governs ranges::empty?}\\[-1mm]{\scriptsize\CornellClauseRef{26.3.12}} & Explains the traversal or view contract for ranges::empty; prove domain validity, lifetime, sentinel compatibility, and any borrowing or invalidation property.\\\hline
\textbf{What contract governs ranges::data?}\\[-1mm]{\scriptsize\CornellClauseRef{26.3.13}} & Explains the traversal or view contract for ranges::data; prove domain validity, lifetime, sentinel compatibility, and any borrowing or invalidation property.\\\hline
\textbf{What contract governs ranges::cdata?}\\[-1mm]{\scriptsize\CornellClauseRef{26.3.14}} & Explains the traversal or view contract for ranges::cdata; prove domain validity, lifetime, sentinel compatibility, and any borrowing or invalidation property.\\\hline
\end{longtable}
\begin{CornellReferenceSummaryBox}Range access supplies the library semantics portion of this clause. The key review move is to connect its local wording to the clause-wide model, then classify each condition by normative force before relying on it.\end{CornellReferenceSummaryBox}
\Needspace{8\baselineskip}
\subsection*{26.4 Range requirements}
\begin{longtable}{@{}>{\RaggedRight\arraybackslash\columncolor{CornellReferenceCueGray}}p{0.28\textwidth}>{\RaggedRight\arraybackslash}p{0.64\textwidth}@{}}
\rowcolor{CornellReferenceNavy}\color{white}\textbf{Cue / recall prompt} & \color{white}\textbf{Notes}\\\endfirsthead
\rowcolor{CornellReferenceNavy}\color{white}\textbf{Cue / recall prompt} & \color{white}\textbf{Notes (continued)}\\\endhead
\arrayrulecolor{CornellReferenceLineGray}
\textbf{What contract governs General?}\\[-1mm]{\scriptsize\CornellClauseRef{26.4.1}} & Frames the terminology, applicability, and relationships needed to interpret General; detailed rules remain in the following subclauses.\\\hline
\textbf{What contract governs Ranges?}\\[-1mm]{\scriptsize\CornellClauseRef{26.4.2}} & Explains the traversal or view contract for Ranges; prove domain validity, lifetime, sentinel compatibility, and any borrowing or invalidation property.\\\hline
\textbf{What contract governs Sized ranges?}\\[-1mm]{\scriptsize\CornellClauseRef{26.4.3}} & Explains the traversal or view contract for Sized ranges; prove domain validity, lifetime, sentinel compatibility, and any borrowing or invalidation property.\\\hline
\textbf{What contract governs Views?}\\[-1mm]{\scriptsize\CornellClauseRef{26.4.4}} & Explains the traversal or view contract for Views; prove domain validity, lifetime, sentinel compatibility, and any borrowing or invalidation property.\\\hline
\textbf{What contract governs Other range refinements?}\\[-1mm]{\scriptsize\CornellClauseRef{26.4.5}} & Explains the traversal or view contract for Other range refinements; prove domain validity, lifetime, sentinel compatibility, and any borrowing or invalidation property.\\\hline
\end{longtable}
\begin{CornellReferenceSummaryBox}Range requirements supplies the requirements portion of this clause. The key review move is to connect its local wording to the clause-wide model, then classify each condition by normative force before relying on it.\end{CornellReferenceSummaryBox}
\Needspace{8\baselineskip}
\subsection*{26.5 Range utilities}
\begin{longtable}{@{}>{\RaggedRight\arraybackslash\columncolor{CornellReferenceCueGray}}p{0.28\textwidth}>{\RaggedRight\arraybackslash}p{0.64\textwidth}@{}}
\rowcolor{CornellReferenceNavy}\color{white}\textbf{Cue / recall prompt} & \color{white}\textbf{Notes}\\\endfirsthead
\rowcolor{CornellReferenceNavy}\color{white}\textbf{Cue / recall prompt} & \color{white}\textbf{Notes (continued)}\\\endhead
\arrayrulecolor{CornellReferenceLineGray}
\textbf{What contract governs General?}\\[-1mm]{\scriptsize\CornellClauseRef{26.5.1}} & Frames the terminology, applicability, and relationships needed to interpret General; detailed rules remain in the following subclauses.\\\hline
\textbf{What contract governs Helper concepts?}\\[-1mm]{\scriptsize\CornellClauseRef{26.5.2}} & Defines or applies the generic requirement Helper concepts; syntactic satisfaction must be paired with every stated semantic and equality-preservation expectation.\\\hline
\textbf{What contract governs View interface?}\\[-1mm]{\scriptsize\CornellClauseRef{26.5.3}} & Explains the traversal or view contract for View interface; prove domain validity, lifetime, sentinel compatibility, and any borrowing or invalidation property.\\\hline
\textbf{What contract governs Sub-ranges?}\\[-1mm]{\scriptsize\CornellClauseRef{26.5.4}} & Explains the traversal or view contract for Sub-ranges; prove domain validity, lifetime, sentinel compatibility, and any borrowing or invalidation property.\\\hline
\textbf{What contract governs Dangling iterator handling?}\\[-1mm]{\scriptsize\CornellClauseRef{26.5.5}} & Explains the traversal or view contract for Dangling iterator handling; prove domain validity, lifetime, sentinel compatibility, and any borrowing or invalidation property.\\\hline
\textbf{What contract governs Class template elements\_\allowbreak{}of?}\\[-1mm]{\scriptsize\CornellClauseRef{26.5.6}} & Specifies the facility family Class template elements\_\allowbreak{}of. Read its declarations together with mandates, preconditions, effects, results, exceptions, complexity, and lifetime rules.\\\hline
\textbf{When is Range conversions permitted?}\\[-1mm]{\scriptsize\CornellClauseRef{26.5.7}} & Explains the traversal or view contract for Range conversions; prove domain validity, lifetime, sentinel compatibility, and any borrowing or invalidation property.\\\hline
\end{longtable}
\begin{CornellReferenceSummaryBox}Range utilities supplies the library semantics portion of this clause. The key review move is to connect its local wording to the clause-wide model, then classify each condition by normative force before relying on it.\end{CornellReferenceSummaryBox}
\Needspace{5\baselineskip}\paragraph{Detailed coverage ledger.}
\begin{description}[style=nextline,leftmargin=2.8cm,font=\normalfont\bfseries\color{CornellReferenceTeal}]
\item[26.5.3.1] \textbf{General.} Frames the terminology, applicability, and relationships needed to interpret General; detailed rules remain in the following subclauses.
\item[26.5.3.2] \textbf{Members.} Specifies the facility family Members. Read its declarations together with mandates, preconditions, effects, results, exceptions, complexity, and lifetime rules.
\item[26.5.4.1] \textbf{General.} Frames the terminology, applicability, and relationships needed to interpret General; detailed rules remain in the following subclauses.
\item[26.5.4.2] \textbf{Constructors and conversions.} Specifies how Constructors and conversions establishes object state, including participation constraints, initialization effects, and exceptional exits.
\item[26.5.4.3] \textbf{Accessors.} Specifies the facility family Accessors. Read its declarations together with mandates, preconditions, effects, results, exceptions, complexity, and lifetime rules.
\item[26.5.7.1] \textbf{General.} Frames the terminology, applicability, and relationships needed to interpret General; detailed rules remain in the following subclauses.
\item[26.5.7.2] \textbf{ranges::to.} Explains the traversal or view contract for ranges::to; prove domain validity, lifetime, sentinel compatibility, and any borrowing or invalidation property.
\item[26.5.7.3] \textbf{ranges::to adaptors.} Explains the traversal or view contract for ranges::to adaptors; prove domain validity, lifetime, sentinel compatibility, and any borrowing or invalidation property.
\end{description}
\Needspace{8\baselineskip}
\subsection*{26.6 Range factories}
\begin{longtable}{@{}>{\RaggedRight\arraybackslash\columncolor{CornellReferenceCueGray}}p{0.28\textwidth}>{\RaggedRight\arraybackslash}p{0.64\textwidth}@{}}
\rowcolor{CornellReferenceNavy}\color{white}\textbf{Cue / recall prompt} & \color{white}\textbf{Notes}\\\endfirsthead
\rowcolor{CornellReferenceNavy}\color{white}\textbf{Cue / recall prompt} & \color{white}\textbf{Notes (continued)}\\\endhead
\arrayrulecolor{CornellReferenceLineGray}
\textbf{What contract governs General?}\\[-1mm]{\scriptsize\CornellClauseRef{26.6.1}} & Frames the terminology, applicability, and relationships needed to interpret General; detailed rules remain in the following subclauses.\\\hline
\textbf{What contract governs Empty view?}\\[-1mm]{\scriptsize\CornellClauseRef{26.6.2}} & Explains the traversal or view contract for Empty view; prove domain validity, lifetime, sentinel compatibility, and any borrowing or invalidation property.\\\hline
\textbf{What contract governs Single view?}\\[-1mm]{\scriptsize\CornellClauseRef{26.6.3}} & Explains the traversal or view contract for Single view; prove domain validity, lifetime, sentinel compatibility, and any borrowing or invalidation property.\\\hline
\textbf{What contract governs Iota view?}\\[-1mm]{\scriptsize\CornellClauseRef{26.6.4}} & Explains the traversal or view contract for Iota view; prove domain validity, lifetime, sentinel compatibility, and any borrowing or invalidation property.\\\hline
\textbf{What contract governs Repeat view?}\\[-1mm]{\scriptsize\CornellClauseRef{26.6.5}} & Explains the traversal or view contract for Repeat view; prove domain validity, lifetime, sentinel compatibility, and any borrowing or invalidation property.\\\hline
\textbf{What contract governs Istream view?}\\[-1mm]{\scriptsize\CornellClauseRef{26.6.6}} & Explains the traversal or view contract for Istream view; prove domain validity, lifetime, sentinel compatibility, and any borrowing or invalidation property.\\\hline
\end{longtable}
\begin{CornellReferenceSummaryBox}Range factories supplies the library semantics portion of this clause. The key review move is to connect its local wording to the clause-wide model, then classify each condition by normative force before relying on it.\end{CornellReferenceSummaryBox}
\Needspace{5\baselineskip}\paragraph{Detailed coverage ledger.}
\begin{description}[style=nextline,leftmargin=2.8cm,font=\normalfont\bfseries\color{CornellReferenceTeal}]
\item[26.6.2.1] \textbf{Overview.} Frames the terminology, applicability, and relationships needed to interpret Overview; detailed rules remain in the following subclauses.
\item[26.6.2.2] \textbf{Class template empty\_\allowbreak{}view.} Explains the traversal or view contract for Class template empty\_\allowbreak{}view; prove domain validity, lifetime, sentinel compatibility, and any borrowing or invalidation property.
\item[26.6.3.1] \textbf{Overview.} Frames the terminology, applicability, and relationships needed to interpret Overview; detailed rules remain in the following subclauses.
\item[26.6.3.2] \textbf{Class template single\_\allowbreak{}view.} Explains the traversal or view contract for Class template single\_\allowbreak{}view; prove domain validity, lifetime, sentinel compatibility, and any borrowing or invalidation property.
\item[26.6.4.1] \textbf{Overview.} Frames the terminology, applicability, and relationships needed to interpret Overview; detailed rules remain in the following subclauses.
\item[26.6.4.2] \textbf{Class template iota\_\allowbreak{}view.} Explains the traversal or view contract for Class template iota\_\allowbreak{}view; prove domain validity, lifetime, sentinel compatibility, and any borrowing or invalidation property.
\item[26.6.4.3] \textbf{Class iota\_\allowbreak{}view::iterator.} Explains the traversal or view contract for Class iota\_\allowbreak{}view::iterator; prove domain validity, lifetime, sentinel compatibility, and any borrowing or invalidation property.
\item[26.6.4.4] \textbf{Class iota\_\allowbreak{}view::sentinel.} Explains the traversal or view contract for Class iota\_\allowbreak{}view::sentinel; prove domain validity, lifetime, sentinel compatibility, and any borrowing or invalidation property.
\item[26.6.5.1] \textbf{Overview.} Frames the terminology, applicability, and relationships needed to interpret Overview; detailed rules remain in the following subclauses.
\item[26.6.5.2] \textbf{Class template repeat\_\allowbreak{}view.} Explains the traversal or view contract for Class template repeat\_\allowbreak{}view; prove domain validity, lifetime, sentinel compatibility, and any borrowing or invalidation property.
\item[26.6.5.3] \textbf{Class repeat\_\allowbreak{}view::iterator.} Explains the traversal or view contract for Class repeat\_\allowbreak{}view::iterator; prove domain validity, lifetime, sentinel compatibility, and any borrowing or invalidation property.
\item[26.6.6.1] \textbf{Overview.} Frames the terminology, applicability, and relationships needed to interpret Overview; detailed rules remain in the following subclauses.
\item[26.6.6.2] \textbf{Class template basic\_\allowbreak{}istream\_\allowbreak{}view.} Explains the traversal or view contract for Class template basic\_\allowbreak{}istream\_\allowbreak{}view; prove domain validity, lifetime, sentinel compatibility, and any borrowing or invalidation property.
\item[26.6.6.3] \textbf{Class basic\_\allowbreak{}istream\_\allowbreak{}view::iterator.} Explains the traversal or view contract for Class basic\_\allowbreak{}istream\_\allowbreak{}view::iterator; prove domain validity, lifetime, sentinel compatibility, and any borrowing or invalidation property.
\end{description}
\Needspace{8\baselineskip}
\subsection*{26.7 Range adaptors}
\begin{longtable}{@{}>{\RaggedRight\arraybackslash\columncolor{CornellReferenceCueGray}}p{0.28\textwidth}>{\RaggedRight\arraybackslash}p{0.64\textwidth}@{}}
\rowcolor{CornellReferenceNavy}\color{white}\textbf{Cue / recall prompt} & \color{white}\textbf{Notes}\\\endfirsthead
\rowcolor{CornellReferenceNavy}\color{white}\textbf{Cue / recall prompt} & \color{white}\textbf{Notes (continued)}\\\endhead
\arrayrulecolor{CornellReferenceLineGray}
\textbf{What contract governs General?}\\[-1mm]{\scriptsize\CornellClauseRef{26.7.1}} & Frames the terminology, applicability, and relationships needed to interpret General; detailed rules remain in the following subclauses.\\\hline
\textbf{What contract governs Range adaptor objects?}\\[-1mm]{\scriptsize\CornellClauseRef{26.7.2}} & Explains the traversal or view contract for Range adaptor objects; prove domain validity, lifetime, sentinel compatibility, and any borrowing or invalidation property.\\\hline
\textbf{What contract governs Movable wrapper?}\\[-1mm]{\scriptsize\CornellClauseRef{26.7.3}} & Specifies the facility family Movable wrapper. Read its declarations together with mandates, preconditions, effects, results, exceptions, complexity, and lifetime rules.\\\hline
\textbf{What contract governs Non-propagating cache?}\\[-1mm]{\scriptsize\CornellClauseRef{26.7.4}} & Specifies the facility family Non-propagating cache. Read its declarations together with mandates, preconditions, effects, results, exceptions, complexity, and lifetime rules.\\\hline
\textbf{What contract governs Range adaptor helpers?}\\[-1mm]{\scriptsize\CornellClauseRef{26.7.5}} & Explains the traversal or view contract for Range adaptor helpers; prove domain validity, lifetime, sentinel compatibility, and any borrowing or invalidation property.\\\hline
\textbf{What contract governs All view?}\\[-1mm]{\scriptsize\CornellClauseRef{26.7.6}} & Explains the traversal or view contract for All view; prove domain validity, lifetime, sentinel compatibility, and any borrowing or invalidation property.\\\hline
\textbf{What contract governs As rvalue view?}\\[-1mm]{\scriptsize\CornellClauseRef{26.7.7}} & Explains the traversal or view contract for As rvalue view; prove domain validity, lifetime, sentinel compatibility, and any borrowing or invalidation property.\\\hline
\textbf{What contract governs Filter view?}\\[-1mm]{\scriptsize\CornellClauseRef{26.7.8}} & Explains the traversal or view contract for Filter view; prove domain validity, lifetime, sentinel compatibility, and any borrowing or invalidation property.\\\hline
\textbf{What contract governs Transform view?}\\[-1mm]{\scriptsize\CornellClauseRef{26.7.9}} & Explains the traversal or view contract for Transform view; prove domain validity, lifetime, sentinel compatibility, and any borrowing or invalidation property.\\\hline
\textbf{What contract governs Take view?}\\[-1mm]{\scriptsize\CornellClauseRef{26.7.10}} & Explains the traversal or view contract for Take view; prove domain validity, lifetime, sentinel compatibility, and any borrowing or invalidation property.\\\hline
\textbf{What contract governs Take while view?}\\[-1mm]{\scriptsize\CornellClauseRef{26.7.11}} & Explains the traversal or view contract for Take while view; prove domain validity, lifetime, sentinel compatibility, and any borrowing or invalidation property.\\\hline
\textbf{What contract governs Drop view?}\\[-1mm]{\scriptsize\CornellClauseRef{26.7.12}} & Explains the traversal or view contract for Drop view; prove domain validity, lifetime, sentinel compatibility, and any borrowing or invalidation property.\\\hline
\textbf{What contract governs Drop while view?}\\[-1mm]{\scriptsize\CornellClauseRef{26.7.13}} & Explains the traversal or view contract for Drop while view; prove domain validity, lifetime, sentinel compatibility, and any borrowing or invalidation property.\\\hline
\textbf{What contract governs Join view?}\\[-1mm]{\scriptsize\CornellClauseRef{26.7.14}} & Explains the traversal or view contract for Join view; prove domain validity, lifetime, sentinel compatibility, and any borrowing or invalidation property.\\\hline
\textbf{What contract governs Join with view?}\\[-1mm]{\scriptsize\CornellClauseRef{26.7.15}} & Explains the traversal or view contract for Join with view; prove domain validity, lifetime, sentinel compatibility, and any borrowing or invalidation property.\\\hline
\textbf{What contract governs Lazy split view?}\\[-1mm]{\scriptsize\CornellClauseRef{26.7.16}} & Explains the traversal or view contract for Lazy split view; prove domain validity, lifetime, sentinel compatibility, and any borrowing or invalidation property.\\\hline
\textbf{What contract governs Split view?}\\[-1mm]{\scriptsize\CornellClauseRef{26.7.17}} & Explains the traversal or view contract for Split view; prove domain validity, lifetime, sentinel compatibility, and any borrowing or invalidation property.\\\hline
\textbf{What contract governs Counted view?}\\[-1mm]{\scriptsize\CornellClauseRef{26.7.18}} & Explains the traversal or view contract for Counted view; prove domain validity, lifetime, sentinel compatibility, and any borrowing or invalidation property.\\\hline
\textbf{What contract governs Common view?}\\[-1mm]{\scriptsize\CornellClauseRef{26.7.19}} & Explains the traversal or view contract for Common view; prove domain validity, lifetime, sentinel compatibility, and any borrowing or invalidation property.\\\hline
\textbf{What contract governs Reverse view?}\\[-1mm]{\scriptsize\CornellClauseRef{26.7.20}} & Explains the traversal or view contract for Reverse view; prove domain validity, lifetime, sentinel compatibility, and any borrowing or invalidation property.\\\hline
\textbf{What contract governs As const view?}\\[-1mm]{\scriptsize\CornellClauseRef{26.7.21}} & Explains the traversal or view contract for As const view; prove domain validity, lifetime, sentinel compatibility, and any borrowing or invalidation property.\\\hline
\textbf{What contract governs Elements view?}\\[-1mm]{\scriptsize\CornellClauseRef{26.7.22}} & Explains the traversal or view contract for Elements view; prove domain validity, lifetime, sentinel compatibility, and any borrowing or invalidation property.\\\hline
\textbf{What contract governs Enumerate view?}\\[-1mm]{\scriptsize\CornellClauseRef{26.7.23}} & Explains the traversal or view contract for Enumerate view; prove domain validity, lifetime, sentinel compatibility, and any borrowing or invalidation property.\\\hline
\textbf{What contract governs Zip view?}\\[-1mm]{\scriptsize\CornellClauseRef{26.7.24}} & Explains the traversal or view contract for Zip view; prove domain validity, lifetime, sentinel compatibility, and any borrowing or invalidation property.\\\hline
\textbf{What contract governs Zip transform view?}\\[-1mm]{\scriptsize\CornellClauseRef{26.7.25}} & Explains the traversal or view contract for Zip transform view; prove domain validity, lifetime, sentinel compatibility, and any borrowing or invalidation property.\\\hline
\textbf{What contract governs Adjacent view?}\\[-1mm]{\scriptsize\CornellClauseRef{26.7.26}} & Explains the traversal or view contract for Adjacent view; prove domain validity, lifetime, sentinel compatibility, and any borrowing or invalidation property.\\\hline
\textbf{What contract governs Adjacent transform view?}\\[-1mm]{\scriptsize\CornellClauseRef{26.7.27}} & Explains the traversal or view contract for Adjacent transform view; prove domain validity, lifetime, sentinel compatibility, and any borrowing or invalidation property.\\\hline
\textbf{What contract governs Chunk view?}\\[-1mm]{\scriptsize\CornellClauseRef{26.7.28}} & Explains the traversal or view contract for Chunk view; prove domain validity, lifetime, sentinel compatibility, and any borrowing or invalidation property.\\\hline
\textbf{What contract governs Slide view?}\\[-1mm]{\scriptsize\CornellClauseRef{26.7.29}} & Explains the traversal or view contract for Slide view; prove domain validity, lifetime, sentinel compatibility, and any borrowing or invalidation property.\\\hline
\textbf{What contract governs Chunk by view?}\\[-1mm]{\scriptsize\CornellClauseRef{26.7.30}} & Explains the traversal or view contract for Chunk by view; prove domain validity, lifetime, sentinel compatibility, and any borrowing or invalidation property.\\\hline
\textbf{What contract governs Stride view?}\\[-1mm]{\scriptsize\CornellClauseRef{26.7.31}} & Explains the traversal or view contract for Stride view; prove domain validity, lifetime, sentinel compatibility, and any borrowing or invalidation property.\\\hline
\textbf{What contract governs Cartesian product view?}\\[-1mm]{\scriptsize\CornellClauseRef{26.7.32}} & Explains the traversal or view contract for Cartesian product view; prove domain validity, lifetime, sentinel compatibility, and any borrowing or invalidation property.\\\hline
\end{longtable}
\begin{CornellReferenceSummaryBox}Range adaptors supplies the library semantics portion of this clause. The key review move is to connect its local wording to the clause-wide model, then classify each condition by normative force before relying on it.\end{CornellReferenceSummaryBox}
\Needspace{5\baselineskip}\paragraph{Detailed coverage ledger.}
\begin{description}[style=nextline,leftmargin=2.8cm,font=\normalfont\bfseries\color{CornellReferenceTeal}]
\item[26.7.6.1] \textbf{General.} Frames the terminology, applicability, and relationships needed to interpret General; detailed rules remain in the following subclauses.
\item[26.7.6.2] \textbf{Class template ref\_\allowbreak{}view.} Explains the traversal or view contract for Class template ref\_\allowbreak{}view; prove domain validity, lifetime, sentinel compatibility, and any borrowing or invalidation property.
\item[26.7.6.3] \textbf{Class template owning\_\allowbreak{}view.} Explains the traversal or view contract for Class template owning\_\allowbreak{}view; prove domain validity, lifetime, sentinel compatibility, and any borrowing or invalidation property.
\item[26.7.7.1] \textbf{Overview.} Frames the terminology, applicability, and relationships needed to interpret Overview; detailed rules remain in the following subclauses.
\item[26.7.7.2] \textbf{Class template as\_\allowbreak{}rvalue\_\allowbreak{}view.} Explains the traversal or view contract for Class template as\_\allowbreak{}rvalue\_\allowbreak{}view; prove domain validity, lifetime, sentinel compatibility, and any borrowing or invalidation property.
\item[26.7.8.1] \textbf{Overview.} Frames the terminology, applicability, and relationships needed to interpret Overview; detailed rules remain in the following subclauses.
\item[26.7.8.2] \textbf{Class template filter\_\allowbreak{}view.} Explains the traversal or view contract for Class template filter\_\allowbreak{}view; prove domain validity, lifetime, sentinel compatibility, and any borrowing or invalidation property.
\item[26.7.8.3] \textbf{Class filter\_\allowbreak{}view::iterator.} Explains the traversal or view contract for Class filter\_\allowbreak{}view::iterator; prove domain validity, lifetime, sentinel compatibility, and any borrowing or invalidation property.
\item[26.7.8.4] \textbf{Class filter\_\allowbreak{}view::sentinel.} Explains the traversal or view contract for Class filter\_\allowbreak{}view::sentinel; prove domain validity, lifetime, sentinel compatibility, and any borrowing or invalidation property.
\item[26.7.9.1] \textbf{Overview.} Frames the terminology, applicability, and relationships needed to interpret Overview; detailed rules remain in the following subclauses.
\item[26.7.9.2] \textbf{Class template transform\_\allowbreak{}view.} Explains the traversal or view contract for Class template transform\_\allowbreak{}view; prove domain validity, lifetime, sentinel compatibility, and any borrowing or invalidation property.
\item[26.7.9.3] \textbf{Class template transform\_\allowbreak{}view::iterator.} Explains the traversal or view contract for Class template transform\_\allowbreak{}view::iterator; prove domain validity, lifetime, sentinel compatibility, and any borrowing or invalidation property.
\item[26.7.9.4] \textbf{Class template transform\_\allowbreak{}view::sentinel.} Explains the traversal or view contract for Class template transform\_\allowbreak{}view::sentinel; prove domain validity, lifetime, sentinel compatibility, and any borrowing or invalidation property.
\item[26.7.10.1] \textbf{Overview.} Frames the terminology, applicability, and relationships needed to interpret Overview; detailed rules remain in the following subclauses.
\item[26.7.10.2] \textbf{Class template take\_\allowbreak{}view.} Explains the traversal or view contract for Class template take\_\allowbreak{}view; prove domain validity, lifetime, sentinel compatibility, and any borrowing or invalidation property.
\item[26.7.10.3] \textbf{Class template take\_\allowbreak{}view::sentinel.} Explains the traversal or view contract for Class template take\_\allowbreak{}view::sentinel; prove domain validity, lifetime, sentinel compatibility, and any borrowing or invalidation property.
\item[26.7.11.1] \textbf{Overview.} Frames the terminology, applicability, and relationships needed to interpret Overview; detailed rules remain in the following subclauses.
\item[26.7.11.2] \textbf{Class template take\_\allowbreak{}while\_\allowbreak{}view.} Explains the traversal or view contract for Class template take\_\allowbreak{}while\_\allowbreak{}view; prove domain validity, lifetime, sentinel compatibility, and any borrowing or invalidation property.
\item[26.7.11.3] \textbf{Class template take\_\allowbreak{}while\_\allowbreak{}view::sentinel.} Explains the traversal or view contract for Class template take\_\allowbreak{}while\_\allowbreak{}view::sentinel; prove domain validity, lifetime, sentinel compatibility, and any borrowing or invalidation property.
\item[26.7.12.1] \textbf{Overview.} Frames the terminology, applicability, and relationships needed to interpret Overview; detailed rules remain in the following subclauses.
\item[26.7.12.2] \textbf{Class template drop\_\allowbreak{}view.} Explains the traversal or view contract for Class template drop\_\allowbreak{}view; prove domain validity, lifetime, sentinel compatibility, and any borrowing or invalidation property.
\item[26.7.13.1] \textbf{Overview.} Frames the terminology, applicability, and relationships needed to interpret Overview; detailed rules remain in the following subclauses.
\item[26.7.13.2] \textbf{Class template drop\_\allowbreak{}while\_\allowbreak{}view.} Explains the traversal or view contract for Class template drop\_\allowbreak{}while\_\allowbreak{}view; prove domain validity, lifetime, sentinel compatibility, and any borrowing or invalidation property.
\item[26.7.14.1] \textbf{Overview.} Frames the terminology, applicability, and relationships needed to interpret Overview; detailed rules remain in the following subclauses.
\item[26.7.14.2] \textbf{Class template join\_\allowbreak{}view.} Explains the traversal or view contract for Class template join\_\allowbreak{}view; prove domain validity, lifetime, sentinel compatibility, and any borrowing or invalidation property.
\item[26.7.14.3] \textbf{Class template join\_\allowbreak{}view::iterator.} Explains the traversal or view contract for Class template join\_\allowbreak{}view::iterator; prove domain validity, lifetime, sentinel compatibility, and any borrowing or invalidation property.
\item[26.7.14.4] \textbf{Class template join\_\allowbreak{}view::sentinel.} Explains the traversal or view contract for Class template join\_\allowbreak{}view::sentinel; prove domain validity, lifetime, sentinel compatibility, and any borrowing or invalidation property.
\item[26.7.15.1] \textbf{Overview.} Frames the terminology, applicability, and relationships needed to interpret Overview; detailed rules remain in the following subclauses.
\item[26.7.15.2] \textbf{Class template join\_\allowbreak{}with\_\allowbreak{}view.} Explains the traversal or view contract for Class template join\_\allowbreak{}with\_\allowbreak{}view; prove domain validity, lifetime, sentinel compatibility, and any borrowing or invalidation property.
\item[26.7.15.3] \textbf{Class template join\_\allowbreak{}with\_\allowbreak{}view::iterator.} Explains the traversal or view contract for Class template join\_\allowbreak{}with\_\allowbreak{}view::iterator; prove domain validity, lifetime, sentinel compatibility, and any borrowing or invalidation property.
\item[26.7.15.4] \textbf{Class template join\_\allowbreak{}with\_\allowbreak{}view::sentinel.} Explains the traversal or view contract for Class template join\_\allowbreak{}with\_\allowbreak{}view::sentinel; prove domain validity, lifetime, sentinel compatibility, and any borrowing or invalidation property.
\item[26.7.16.1] \textbf{Overview.} Frames the terminology, applicability, and relationships needed to interpret Overview; detailed rules remain in the following subclauses.
\item[26.7.16.2] \textbf{Class template lazy\_\allowbreak{}split\_\allowbreak{}view.} Explains the traversal or view contract for Class template lazy\_\allowbreak{}split\_\allowbreak{}view; prove domain validity, lifetime, sentinel compatibility, and any borrowing or invalidation property.
\item[26.7.16.3] \textbf{Class template lazy\_\allowbreak{}split\_\allowbreak{}view::outer-iterator.} Explains the traversal or view contract for Class template lazy\_\allowbreak{}split\_\allowbreak{}view::outer-iterator; prove domain validity, lifetime, sentinel compatibility, and any borrowing or invalidation property.
\item[26.7.16.4] \textbf{Class lazy\_\allowbreak{}split\_\allowbreak{}view::outer-iterator::value\_\allowbreak{}type.} Explains the traversal or view contract for Class lazy\_\allowbreak{}split\_\allowbreak{}view::outer-iterator::value\_\allowbreak{}type; prove domain validity, lifetime, sentinel compatibility, and any borrowing or invalidation property.
\item[26.7.16.5] \textbf{Class template lazy\_\allowbreak{}split\_\allowbreak{}view::inner-iterator.} Explains the traversal or view contract for Class template lazy\_\allowbreak{}split\_\allowbreak{}view::inner-iterator; prove domain validity, lifetime, sentinel compatibility, and any borrowing or invalidation property.
\item[26.7.17.1] \textbf{Overview.} Frames the terminology, applicability, and relationships needed to interpret Overview; detailed rules remain in the following subclauses.
\item[26.7.17.2] \textbf{Class template split\_\allowbreak{}view.} Explains the traversal or view contract for Class template split\_\allowbreak{}view; prove domain validity, lifetime, sentinel compatibility, and any borrowing or invalidation property.
\item[26.7.17.3] \textbf{Class split\_\allowbreak{}view::iterator.} Explains the traversal or view contract for Class split\_\allowbreak{}view::iterator; prove domain validity, lifetime, sentinel compatibility, and any borrowing or invalidation property.
\item[26.7.17.4] \textbf{Class split\_\allowbreak{}view::sentinel.} Explains the traversal or view contract for Class split\_\allowbreak{}view::sentinel; prove domain validity, lifetime, sentinel compatibility, and any borrowing or invalidation property.
\item[26.7.19.1] \textbf{Overview.} Frames the terminology, applicability, and relationships needed to interpret Overview; detailed rules remain in the following subclauses.
\item[26.7.19.2] \textbf{Class template common\_\allowbreak{}view.} Explains the traversal or view contract for Class template common\_\allowbreak{}view; prove domain validity, lifetime, sentinel compatibility, and any borrowing or invalidation property.
\item[26.7.20.1] \textbf{Overview.} Frames the terminology, applicability, and relationships needed to interpret Overview; detailed rules remain in the following subclauses.
\item[26.7.20.2] \textbf{Class template reverse\_\allowbreak{}view.} Explains the traversal or view contract for Class template reverse\_\allowbreak{}view; prove domain validity, lifetime, sentinel compatibility, and any borrowing or invalidation property.
\item[26.7.21.1] \textbf{Overview.} Frames the terminology, applicability, and relationships needed to interpret Overview; detailed rules remain in the following subclauses.
\item[26.7.21.2] \textbf{Class template as\_\allowbreak{}const\_\allowbreak{}view.} Explains the traversal or view contract for Class template as\_\allowbreak{}const\_\allowbreak{}view; prove domain validity, lifetime, sentinel compatibility, and any borrowing or invalidation property.
\item[26.7.22.1] \textbf{Overview.} Frames the terminology, applicability, and relationships needed to interpret Overview; detailed rules remain in the following subclauses.
\item[26.7.22.2] \textbf{Class template elements\_\allowbreak{}view.} Explains the traversal or view contract for Class template elements\_\allowbreak{}view; prove domain validity, lifetime, sentinel compatibility, and any borrowing or invalidation property.
\item[26.7.22.3] \textbf{Class template elements\_\allowbreak{}view::iterator.} Explains the traversal or view contract for Class template elements\_\allowbreak{}view::iterator; prove domain validity, lifetime, sentinel compatibility, and any borrowing or invalidation property.
\item[26.7.22.4] \textbf{Class template elements\_\allowbreak{}view::sentinel.} Explains the traversal or view contract for Class template elements\_\allowbreak{}view::sentinel; prove domain validity, lifetime, sentinel compatibility, and any borrowing or invalidation property.
\item[26.7.23.1] \textbf{Overview.} Frames the terminology, applicability, and relationships needed to interpret Overview; detailed rules remain in the following subclauses.
\item[26.7.23.2] \textbf{Class template enumerate\_\allowbreak{}view.} Explains the traversal or view contract for Class template enumerate\_\allowbreak{}view; prove domain validity, lifetime, sentinel compatibility, and any borrowing or invalidation property.
\item[26.7.23.3] \textbf{Class template enumerate\_\allowbreak{}view::iterator.} Explains the traversal or view contract for Class template enumerate\_\allowbreak{}view::iterator; prove domain validity, lifetime, sentinel compatibility, and any borrowing or invalidation property.
\item[26.7.23.4] \textbf{Class template enumerate\_\allowbreak{}view::sentinel.} Explains the traversal or view contract for Class template enumerate\_\allowbreak{}view::sentinel; prove domain validity, lifetime, sentinel compatibility, and any borrowing or invalidation property.
\item[26.7.24.1] \textbf{Overview.} Frames the terminology, applicability, and relationships needed to interpret Overview; detailed rules remain in the following subclauses.
\item[26.7.24.2] \textbf{Class template zip\_\allowbreak{}view.} Explains the traversal or view contract for Class template zip\_\allowbreak{}view; prove domain validity, lifetime, sentinel compatibility, and any borrowing or invalidation property.
\item[26.7.24.3] \textbf{Class template zip\_\allowbreak{}view::iterator.} Explains the traversal or view contract for Class template zip\_\allowbreak{}view::iterator; prove domain validity, lifetime, sentinel compatibility, and any borrowing or invalidation property.
\item[26.7.24.4] \textbf{Class template zip\_\allowbreak{}view::sentinel.} Explains the traversal or view contract for Class template zip\_\allowbreak{}view::sentinel; prove domain validity, lifetime, sentinel compatibility, and any borrowing or invalidation property.
\item[26.7.25.1] \textbf{Overview.} Frames the terminology, applicability, and relationships needed to interpret Overview; detailed rules remain in the following subclauses.
\item[26.7.25.2] \textbf{Class template zip\_\allowbreak{}transform\_\allowbreak{}view.} Explains the traversal or view contract for Class template zip\_\allowbreak{}transform\_\allowbreak{}view; prove domain validity, lifetime, sentinel compatibility, and any borrowing or invalidation property.
\item[26.7.25.3] \textbf{Class template zip\_\allowbreak{}transform\_\allowbreak{}view::iterator.} Explains the traversal or view contract for Class template zip\_\allowbreak{}transform\_\allowbreak{}view::iterator; prove domain validity, lifetime, sentinel compatibility, and any borrowing or invalidation property.
\item[26.7.25.4] \textbf{Class template zip\_\allowbreak{}transform\_\allowbreak{}view::sentinel.} Explains the traversal or view contract for Class template zip\_\allowbreak{}transform\_\allowbreak{}view::sentinel; prove domain validity, lifetime, sentinel compatibility, and any borrowing or invalidation property.
\item[26.7.26.1] \textbf{Overview.} Frames the terminology, applicability, and relationships needed to interpret Overview; detailed rules remain in the following subclauses.
\item[26.7.26.2] \textbf{Class template adjacent\_\allowbreak{}view.} Explains the traversal or view contract for Class template adjacent\_\allowbreak{}view; prove domain validity, lifetime, sentinel compatibility, and any borrowing or invalidation property.
\item[26.7.26.3] \textbf{Class template adjacent\_\allowbreak{}view::iterator.} Explains the traversal or view contract for Class template adjacent\_\allowbreak{}view::iterator; prove domain validity, lifetime, sentinel compatibility, and any borrowing or invalidation property.
\item[26.7.26.4] \textbf{Class template adjacent\_\allowbreak{}view::sentinel.} Explains the traversal or view contract for Class template adjacent\_\allowbreak{}view::sentinel; prove domain validity, lifetime, sentinel compatibility, and any borrowing or invalidation property.
\item[26.7.27.1] \textbf{Overview.} Frames the terminology, applicability, and relationships needed to interpret Overview; detailed rules remain in the following subclauses.
\item[26.7.27.2] \textbf{Class template adjacent\_\allowbreak{}transform\_\allowbreak{}view.} Explains the traversal or view contract for Class template adjacent\_\allowbreak{}transform\_\allowbreak{}view; prove domain validity, lifetime, sentinel compatibility, and any borrowing or invalidation property.
\item[26.7.27.3] \textbf{Class template adjacent\_\allowbreak{}transform\_\allowbreak{}view::iterator.} Explains the traversal or view contract for Class template adjacent\_\allowbreak{}transform\_\allowbreak{}view::iterator; prove domain validity, lifetime, sentinel compatibility, and any borrowing or invalidation property.
\item[26.7.27.4] \textbf{Class template adjacent\_\allowbreak{}transform\_\allowbreak{}view::sentinel.} Explains the traversal or view contract for Class template adjacent\_\allowbreak{}transform\_\allowbreak{}view::sentinel; prove domain validity, lifetime, sentinel compatibility, and any borrowing or invalidation property.
\item[26.7.28.1] \textbf{Overview.} Frames the terminology, applicability, and relationships needed to interpret Overview; detailed rules remain in the following subclauses.
\item[26.7.28.2] \textbf{Class template chunk\_\allowbreak{}view for input ranges.} Explains the traversal or view contract for Class template chunk\_\allowbreak{}view for input ranges; prove domain validity, lifetime, sentinel compatibility, and any borrowing or invalidation property.
\item[26.7.28.3] \textbf{Class chunk\_\allowbreak{}view::outer-iterator.} Explains the traversal or view contract for Class chunk\_\allowbreak{}view::outer-iterator; prove domain validity, lifetime, sentinel compatibility, and any borrowing or invalidation property.
\item[26.7.28.4] \textbf{Class chunk\_\allowbreak{}view::outer-iterator::value\_\allowbreak{}type.} Explains the traversal or view contract for Class chunk\_\allowbreak{}view::outer-iterator::value\_\allowbreak{}type; prove domain validity, lifetime, sentinel compatibility, and any borrowing or invalidation property.
\item[26.7.28.5] \textbf{Class chunk\_\allowbreak{}view::inner-iterator.} Explains the traversal or view contract for Class chunk\_\allowbreak{}view::inner-iterator; prove domain validity, lifetime, sentinel compatibility, and any borrowing or invalidation property.
\item[26.7.28.6] \textbf{Class template chunk\_\allowbreak{}view for forward ranges.} Explains the traversal or view contract for Class template chunk\_\allowbreak{}view for forward ranges; prove domain validity, lifetime, sentinel compatibility, and any borrowing or invalidation property.
\item[26.7.28.7] \textbf{Class template chunk\_\allowbreak{}view::iterator for forward ranges.} Explains the traversal or view contract for Class template chunk\_\allowbreak{}view::iterator for forward ranges; prove domain validity, lifetime, sentinel compatibility, and any borrowing or invalidation property.
\item[26.7.29.1] \textbf{Overview.} Frames the terminology, applicability, and relationships needed to interpret Overview; detailed rules remain in the following subclauses.
\item[26.7.29.2] \textbf{Class template slide\_\allowbreak{}view.} Explains the traversal or view contract for Class template slide\_\allowbreak{}view; prove domain validity, lifetime, sentinel compatibility, and any borrowing or invalidation property.
\item[26.7.29.3] \textbf{Class template slide\_\allowbreak{}view::iterator.} Explains the traversal or view contract for Class template slide\_\allowbreak{}view::iterator; prove domain validity, lifetime, sentinel compatibility, and any borrowing or invalidation property.
\item[26.7.29.4] \textbf{Class slide\_\allowbreak{}view::sentinel.} Explains the traversal or view contract for Class slide\_\allowbreak{}view::sentinel; prove domain validity, lifetime, sentinel compatibility, and any borrowing or invalidation property.
\item[26.7.30.1] \textbf{Overview.} Frames the terminology, applicability, and relationships needed to interpret Overview; detailed rules remain in the following subclauses.
\item[26.7.30.2] \textbf{Class template chunk\_\allowbreak{}by\_\allowbreak{}view.} Explains the traversal or view contract for Class template chunk\_\allowbreak{}by\_\allowbreak{}view; prove domain validity, lifetime, sentinel compatibility, and any borrowing or invalidation property.
\item[26.7.30.3] \textbf{Class chunk\_\allowbreak{}by\_\allowbreak{}view::iterator.} Explains the traversal or view contract for Class chunk\_\allowbreak{}by\_\allowbreak{}view::iterator; prove domain validity, lifetime, sentinel compatibility, and any borrowing or invalidation property.
\item[26.7.31.1] \textbf{Overview.} Frames the terminology, applicability, and relationships needed to interpret Overview; detailed rules remain in the following subclauses.
\item[26.7.31.2] \textbf{Class template stride\_\allowbreak{}view.} Explains the traversal or view contract for Class template stride\_\allowbreak{}view; prove domain validity, lifetime, sentinel compatibility, and any borrowing or invalidation property.
\item[26.7.31.3] \textbf{Class template stride\_\allowbreak{}view::iterator.} Explains the traversal or view contract for Class template stride\_\allowbreak{}view::iterator; prove domain validity, lifetime, sentinel compatibility, and any borrowing or invalidation property.
\item[26.7.32.1] \textbf{Overview.} Frames the terminology, applicability, and relationships needed to interpret Overview; detailed rules remain in the following subclauses.
\item[26.7.32.2] \textbf{Class template cartesian\_\allowbreak{}product\_\allowbreak{}view.} Explains the traversal or view contract for Class template cartesian\_\allowbreak{}product\_\allowbreak{}view; prove domain validity, lifetime, sentinel compatibility, and any borrowing or invalidation property.
\item[26.7.32.3] \textbf{Class template cartesian\_\allowbreak{}product\_\allowbreak{}view::iterator.} Explains the traversal or view contract for Class template cartesian\_\allowbreak{}product\_\allowbreak{}view::iterator; prove domain validity, lifetime, sentinel compatibility, and any borrowing or invalidation property.
\end{description}
\Needspace{8\baselineskip}
\subsection*{26.8 Range generators}
\begin{longtable}{@{}>{\RaggedRight\arraybackslash\columncolor{CornellReferenceCueGray}}p{0.28\textwidth}>{\RaggedRight\arraybackslash}p{0.64\textwidth}@{}}
\rowcolor{CornellReferenceNavy}\color{white}\textbf{Cue / recall prompt} & \color{white}\textbf{Notes}\\\endfirsthead
\rowcolor{CornellReferenceNavy}\color{white}\textbf{Cue / recall prompt} & \color{white}\textbf{Notes (continued)}\\\endhead
\arrayrulecolor{CornellReferenceLineGray}
\textbf{What contract governs Overview?}\\[-1mm]{\scriptsize\CornellClauseRef{26.8.1}} & Frames the terminology, applicability, and relationships needed to interpret Overview; detailed rules remain in the following subclauses.\\\hline
\textbf{What contract governs Header <generator> synopsis?}\\[-1mm]{\scriptsize\CornellClauseRef{26.8.2}} & Catalogs the declarations made available by Header <generator> synopsis. The synopsis is an interface map; the detailed specifications supply constraints and behavior.\\\hline
\textbf{What contract governs Class template generator?}\\[-1mm]{\scriptsize\CornellClauseRef{26.8.3}} & Specifies the facility family Class template generator. Read its declarations together with mandates, preconditions, effects, results, exceptions, complexity, and lifetime rules.\\\hline
\textbf{What contract governs Members?}\\[-1mm]{\scriptsize\CornellClauseRef{26.8.4}} & Specifies the facility family Members. Read its declarations together with mandates, preconditions, effects, results, exceptions, complexity, and lifetime rules.\\\hline
\textbf{What contract governs Class generator::promise\_\allowbreak{}type?}\\[-1mm]{\scriptsize\CornellClauseRef{26.8.5}} & Defines asynchronous shared-state behavior for Class generator::promise\_\allowbreak{}type; establish producer/consumer ownership, readiness, blocking, exception transport, and destruction effects.\\\hline
\textbf{What contract governs Class generator::iterator?}\\[-1mm]{\scriptsize\CornellClauseRef{26.8.6}} & Explains the traversal or view contract for Class generator::iterator; prove domain validity, lifetime, sentinel compatibility, and any borrowing or invalidation property.\\\hline
\end{longtable}
\begin{CornellReferenceSummaryBox}Range generators supplies the library semantics portion of this clause. The key review move is to connect its local wording to the clause-wide model, then classify each condition by normative force before relying on it.\end{CornellReferenceSummaryBox}
\section{Language-Lawyer and Library-Usage Pitfalls}
\begin{CornellReferencePitfallBox}\begin{itemize}[label=\textcolor{CornellReferenceAmber}{$\blacktriangleright$}]
\item Reading a synopsis as the complete behavioral contract.
\item Ignoring mandates, preconditions, exception behavior, or complexity requirements.
\item Using an exposition-only name as though it were public API.
\item Assuming invalidation, ownership, or thread-safety properties that are not stated.
\item Treating successful compilation as proof that semantic requirements and preconditions are met.
\end{itemize}\end{CornellReferencePitfallBox}
\section{Programmer and Implementation Checklist}
\begin{CornellReferenceChecklistBox}\begin{itemize}[label=$\square$]
\item Identify the declaring header or module exposure for each facility.
\item Check template arguments and mandates before evaluating an operation.
\item Establish every precondition at the call site.
\item Record effects, returned value, postconditions, and exceptions separately.
\item Audit iterator, reference, pointer, and object lifetime after mutation.
\item Verify stated complexity and synchronization assumptions.
\item For \CornellClauseRef{26.1}, verify general against the precise rule category and all cited dependencies.
\item For \CornellClauseRef{26.2}, verify header <ranges> synopsis against the precise rule category and all cited dependencies.
\item For \CornellClauseRef{26.3}, verify range access against the precise rule category and all cited dependencies.
\item For \CornellClauseRef{26.4}, verify range requirements against the precise rule category and all cited dependencies.
\end{itemize}\end{CornellReferenceChecklistBox}
\section{Clause Synthesis}
\begin{CornellReferenceOrientationBox}Defines range concepts, views, adaptors, factories, customization points, and machinery for composing lazy sequence transformations while tracking borrowing, commonality, sizing, and iterator validity. The composable sequence-view layer used by range-aware algorithms. A sound review distinguishes syntax and participation from runtime preconditions, keeps notes and examples non-mandatory, and follows cross-clause dependencies before making a portability claim.\end{CornellReferenceOrientationBox}
\section{Self-Test Questions}
\begin{enumerate}
\item What role does General play in the clause's overall model?
\item Which normative distinctions must be kept separate when applying Header <ranges> synopsis?
\item How would you audit a program or implementation that depends on Range access?
\item Compare Range requirements with Range utilities; what different question does each answer?
\item What role does Range utilities play in the clause's overall model?
\item Which normative distinctions must be kept separate when applying Range factories?
\item How would you audit a program or implementation that depends on Range adaptors?
\item Compare Range generators with General; what different question does each answer?
\item What role does General play in the clause's overall model?
\item Which normative distinctions must be kept separate when applying ranges::begin?
\item How would you audit a program or implementation that depends on ranges::end?
\item Compare ranges::cbegin with ranges::cend; what different question does each answer?
\item What role does ranges::cend play in the clause's overall model?
\item Which normative distinctions must be kept separate when applying ranges::rbegin?
\item Scenario: a program relies on an unstated implementation behavior while using ranges::rend. What must be checked before calling the program portable?
\item Scenario: a program relies on an unstated implementation behavior while using ranges::crbegin. What must be checked before calling the program portable?
\end{enumerate}
\clearpage\section{Answer Key}
\begin{enumerate}
\item General is one part of the clause's library semantics model. Its role is understood by connecting the local rule in 26.7.6.1 to the clause orientation and its dependent concepts.
\item Keep formation and well-formedness separate from runtime preconditions and effects. Also distinguish mandatory wording from notes, implementation choice from undefined behavior, and interface declarations from detailed contracts; see 26.2.
\item Audit Range access by locating the controlling wording in 26.3, listing inputs and type constraints, proving any preconditions, recording effects and failure behavior, and checking lifetime, invalidation, complexity, and synchronization where applicable.
\item Range requirements and Range utilities occupy different steps or facility groups. Analyze each under its own subclause, then follow the explicit dependency between them rather than transferring one set of guarantees to the other; begin at 26.4.
\item Range utilities is one part of the clause's library semantics model. Its role is understood by connecting the local rule in 26.5 to the clause orientation and its dependent concepts.
\item Keep formation and well-formedness separate from runtime preconditions and effects. Also distinguish mandatory wording from notes, implementation choice from undefined behavior, and interface declarations from detailed contracts; see 26.6.
\item Audit Range adaptors by locating the controlling wording in 26.7, listing inputs and type constraints, proving any preconditions, recording effects and failure behavior, and checking lifetime, invalidation, complexity, and synchronization where applicable.
\item Range generators and General occupy different steps or facility groups. Analyze each under its own subclause, then follow the explicit dependency between them rather than transferring one set of guarantees to the other; begin at 26.8.
\item General is one part of the clause's library semantics model. Its role is understood by connecting the local rule in 26.7.6.1 to the clause orientation and its dependent concepts.
\item Keep formation and well-formedness separate from runtime preconditions and effects. Also distinguish mandatory wording from notes, implementation choice from undefined behavior, and interface declarations from detailed contracts; see 26.3.2.
\item Audit ranges::end by locating the controlling wording in 26.3.3, listing inputs and type constraints, proving any preconditions, recording effects and failure behavior, and checking lifetime, invalidation, complexity, and synchronization where applicable.
\item ranges::cbegin and ranges::cend occupy different steps or facility groups. Analyze each under its own subclause, then follow the explicit dependency between them rather than transferring one set of guarantees to the other; begin at 26.3.4.
\item ranges::cend is one part of the clause's library semantics model. Its role is understood by connecting the local rule in 26.3.5 to the clause orientation and its dependent concepts.
\item Keep formation and well-formedness separate from runtime preconditions and effects. Also distinguish mandatory wording from notes, implementation choice from undefined behavior, and interface declarations from detailed contracts; see 26.3.6.
\item First identify the exact requirement in 26.3.7, then classify the relied-on behavior as required, unspecified, implementation-defined, conditionally supported, or undefined. Portability requires satisfying the program obligations and avoiding dependence on a choice not guaranteed in the target environments.
\item First identify the exact requirement in 26.3.8, then classify the relied-on behavior as required, unspecified, implementation-defined, conditionally supported, or undefined. Portability requires satisfying the program obligations and avoiding dependence on a choice not guaranteed in the target environments.
\end{enumerate}
\section{Key-Term Recap}
\begin{longtable}{@{}>{\RaggedRight\arraybackslash}p{0.28\textwidth}>{\RaggedRight\arraybackslash}p{0.64\textwidth}@{}}\toprule\textbf{Term} & \textbf{Working meaning}\\\midrule\endfirsthead\toprule\textbf{Term} & \textbf{Working meaning (continued)}\\\midrule\endhead\bottomrule\endfoot
range & A type for which begin and end expressions provide an iteration domain.\\
view & A lightweight range with specified construction and assignment properties.\\
borrowed\_\allowbreak{}range & A range whose iterators can remain valid after the range object is destroyed.\\
viewable\_\allowbreak{}range & A range suitable for conversion into or use by a view adaptor.\\
range adaptor closure & A callable adaptor object supporting pipeline composition.\\
subrange & A view storing an iterator and sentinel, optionally with a size.\\
\end{longtable}
\CornellTakeaway{Ranges compose safely when view lifetime, borrowing, sizing, and adaptor constraints remain visible through the pipeline.}
\end{document}
